\documentclass[../../main.tex]{subfiles}
\begin{document}

{\bf [解答]}

$\{a_n\}$ を小さい順にならびかえた数列 $\{b_n\}$ とする．

\paragraph{(i)}

$b_1$ の正負により場合分けする．$b_1 \ge 0$ の時は，$b_n$ が単調増加であることから (イ) が成立する．そこで，$b_1 < 0$ の時，(ロ) が成立することを示す．

題意を背理法で示す．$b_n > 0$ と仮定する．
すると $b_k \le 0 < b_{k+1}$ をみたす $k (1 \le k \le n-1)$ が存在する．
\begin{align}
  b_1 < b_2 < \cdots < b_k \le 0 < b_{k+1} < \cdots < b_n
\end{align}
この時，$b_1$ と $b_n$ をとりだしてきて，$b_1-b_n$ 又は $b_n-b_1$ のいずれかが $S$ の要素だが，$b_1 < 0 < b_n$ より
\begin{align}
  b_1 - b_n < b_1 \quad, \quad b_n - b_1 > b_n
\end{align}
となり，いずれも$[b_1,b_n]$の範囲に入らない．これは$b_n$ が単調増加で$b_1$が最小，$b_n$が最大であることに反して矛盾する．以上から $b_n \le 0$ となり，したがって (ロ) が成立する．

以上から題意は示された．

\paragraph{(ii)}

(ロ)が成り立つ時，全ての要素を$-1$倍すれば(イ)のケースに帰着するから，(イ)の場合で証明すれば良い．この時，$S$ の全要素が $0$ 以上だから，$b_k, b_j$ をとってきた時 ($j<k$)，
$b_j - b_k < 0$ は $S$ の要素となりえず，$b_k - b_j$ が $S$ の要素となる．
したがって
\begin{align}
  b_2 - b_1 < b_3 - b_1 < \cdots < b_n - b_1 \ (n-1\text{個}) \label{eq:1}
\end{align}
が全て $S$ の要素となる．

\subparagraph{$1^{\circ} \ b_1 > 0$ の時}

\cref{eq:1}から
\begin{align}
  b_1 \le b_2 - b_1 < b_3 - b_1 < \cdots < b_n - b_1 < b_n
\end{align}
となり，
\begin{align}
  b_2 - b_1 = b_1, \cdots, b_n - b_1 = b_{n-1}
\end{align}
となるから，順に
\begin{align}
  b_2 = 2b_1, \ b_3 = 3b_1, \ \cdots \ b_n = nb_1
\end{align}
となり，たしかに数列 $\{b_n\}$ は初項 $b_1(>0)$，公差 $b_1$ の等差数列である．

\subparagraph{$2^{\circ} \ b_1 = 0$ の時}

$b_2 > b_1 = 0$ だから\cref{eq:1}と同じようにして
\begin{align}
  (0 < ) b_3 - b_2 < b_4 - b_2 < \cdots < b_n - b_2 (< b_n)
\end{align}
が全て $S$ の要素で，
\begin{align}
  b_3 - b_2 = b_2, \cdots, \ b_n - b_2 = b_{n-1}
\end{align}
となり，$2 \le i$ の時 $b_i = i b_2$，$b_1=0$ だから $\{b_n\}$ は初項 $0$，公差 $b_2$ の等差数列である．
\casesend

以上からいずれの場合も $\{b_n\}$ は等差数列である．$\{b_n\}$ は $\{a_n\}$ をならびかえたものだから題意は示された．

\end{document}
